%% =================================================================
%% Beber, Lamon, Saveriano, Fontanelli, Palopoli.
%% "Autonomous Robotic Tissue Palpation and Abnormalities
%%  Characterisation via Ergodic Exploration."
%% IEEE Robotics and Automation Letters, preprint accepted Feb. 2026.
%% DOI: 10.1109/LRA.2026.3673907
%%
%% Reconstruction of page 4 as study notes.
%% Equations (4), (5), (6); plus unnumbered HEDAC smoothing equation.
%% Narrative flows: (continued ergodic control) -> practical impl ->
%%   D. Traj Exec -> 1) Online Elasticity -> 2) Filter Design ->
%%   E. Target Dist Update -> F. Search Termination -> IV. Sim Results
%% =================================================================

\documentclass[11pt]{article}

\usepackage[margin=1in]{geometry}
\usepackage{amsmath, amssymb}
\usepackage{mathrsfs}
\usepackage{bm}
\usepackage{xcolor}
\usepackage{fancyhdr}
\usepackage{url}

\pagestyle{fancy}
\fancyhf{}
\lhead{\small L.~Beber et al.}
\rhead{\small IEEE RA-L Preprint, Accepted Feb.\ 2026}
\cfoot{\small 4 $\to$ \arabic{page}}
\renewcommand{\headrulewidth}{0.4pt}

\setcounter{page}{1}
\setcounter{equation}{3}   % page 3 ended with eq. (3); first numbered eq here is (4)
\setlength{\parskip}{4pt}

\begin{document}

% ---------- Preprint notice strip ----------
\begin{center}
{\footnotesize\itshape
This article has been accepted for publication in IEEE Robotics and
Automation Letters. This is the author's version which has not been
fully edited and content may change prior to final publication.
Citation information: DOI 10.1109/LRA.2026.3673907\par}
\end{center}

\vspace{0.4em}
\hrule
\vspace{0.6em}

\noindent
{\footnotesize\scshape 4 \hfill IEEE Robotics and Automation Letters.\ Preprint Version.\ Accepted February, 2026}

\vspace{1.2em}

% ---------- Continuation of C.1 Ergodic Control from page 3 ----------
\noindent
where $\varphi_{\sigma}$ is a Gaussian kernel. The coverage error with
respect to a target distribution $m(x)$ is defined as the spatial
diffusion of $m(x)$ through the RBF kernel compared with the coverage
density
\begin{equation}
e(x,t) \;=\; (\varphi_{\sigma} * m)(x) \;-\; c(x,t),
\end{equation}
To avoid local minima and improve spatial coordination, HEDAC smooths
this error using a stationary heat equation:
\[
\varrho\,\Delta u(x,t) \;=\; \beta\,u(x,t) + \gamma\,a(x,t) - s(x,t),
\qquad
\dot{z} \;=\; v_{a}\cdot\frac{\nabla u}{\|\nabla u\|},
\]
where $s$ highlights under-sampled regions, and $a$ ensures collision
avoidance between the agents.

\vspace{0.3em}

\noindent\textit{\hspace{1em}2)~Practical Implementation:} Different
from [19], this work implemented the non-stationary diffusion
equation, as was done in [26], for a more locally consistent
exploration behaviour. This means that $\dot{u} \neq 0$ is defined as:
\begin{equation}
\dot{u}(x,\tau) \;=\; \Delta u(x,\tau),
\end{equation}
allowing us to control the desired smoothness. In this formulation,
an initial condition is required: $u(x,0) = s(x,0)$. During the
experiments, a second-order agent is used to generate the trajectory.
The trajectory is updated with a frequency $\Delta T$ chosen to ensure
smooth motion and adequate control loop resolution (typically 100~Hz
in our implementation, 5 times slower than the low-level controller).
The planning phase respects constraints on maximum velocity to ensure
physical plausibility during data collection imposed by the
second-order dynamics inside the ergodic controller.

\vspace{0.6em}

% ---------- D. Trajectory Execution and Data Collection ----------
\noindent\textit{D.~Trajectory Execution and Data Collection}

\vspace{0.2em}

\noindent
While the trajectory is executed, tissue elasticity is estimated
continuously using an Extended Kalman Filter (EKF). Data are sampled
at a rate that depends on the robot motion and the indenter geometry.
Specifically, higher velocities or smaller indenters require denser
sampling to avoid aliasing effects in the stiffness map.

\vspace{0.3em}

\noindent\textit{\hspace{1em}1)~Online Elasticity and Viscosity
Estimation:} Accurately estimating the mechanical response of soft
tissues in real time is a challenging task, as it requires modelling
the contact interaction between the probe and the tissue surface and
inferring tissue parameters from force measurements and robot motion.
In this work, we adopt a viscoelastic contact model to improve the
fidelity of force prediction during dynamic palpation, while
subsequent abnormality detection focuses on the estimated elastic
stiffness To achieve this, we adopt DRM [27], which allows mapping
the three-dimensional contact mechanics into an equivalent
two-dimensional representation. Under this framework, the force
exerted by the tissue during indentation is a nonlinear function of
penetration depth and velocity:
\begin{equation}
F_{\mathrm{TOT}}(d,\dot{d})
  \;=\; \frac{4}{3}\,\frac{E_{f}}{1-\nu^{2}}\,\sqrt{R\,d}\;d
     \;+\; \frac{4\,\eta}{1-\nu}\,\sqrt{R\,d}\;\dot{d},
\end{equation}
where $d$ is the penetration depth, $\dot{d}$ its derivative
(velocity), $E_{f}$ is the Young's modulus of the material, $\eta$ is
the viscosity coefficient, $\nu$ is the Poisson's ratio, and $R$ is
the radius of the spherical indenter [18]. The first term models
elastic deformation according to Hertzian contact mechanics, while
the second term accounts for viscoelastic damping.

\vspace{0.3em}

\noindent\textit{\hspace{1em}2)~Filter Design:} The main challenge
lies in the fact that $d$ is not directly measurable during robotic
interaction. Therefore, we use an EKF to estimate the system state,
which includes both $d$ and $\dot{d}$, as well as the unknown tissue
parameters. The system is modelled using a discrete-time nonlinear
state-space representation, as detailed in [18]. The state vector is
defined as $[d,\,k,\,\lambda,\,\dot{d},\,\dot{k},\,\dot{\lambda},\,
\ddot{k},\,\ddot{\lambda}]^{\top}$, where $\kappa$ and $\lambda$
correspond to the local stiffness and damping parameters,
respectively. The filter assumes as input the measured contact force
and comprises the effective interaction mass. Process and measurement
noise are also included to account for model uncertainties and sensor
imperfections. This filtering approach allows real-time estimation of
the viscoelastic parameters even during continuous contacts with
indentation speeds up to 5~mm/s [28], being suitable for dynamic
exploration tasks.

\vspace{0.6em}

% ---------- E. Target Distribution Update ----------
\noindent\textit{E.~Target Distribution Update}

\vspace{0.2em}

\noindent
Once a sufficient number of new samples are acquired, the GPR model
is updated to refine the target stiffness distribution. Since GPR
inference has cubic time complexity with respect to the number of
samples, we avoid updating the model at every time step. Instead,
updates are performed at a fixed frequency (e.g., 1~Hz), allowing
batch accumulation of new data while keeping computation time
bounded.

\vspace{0.6em}

% ---------- F. Search Termination Criterion ----------
\noindent\textit{F.~Search Termination Criterion}

\vspace{0.2em}

\noindent
Different from existing approaches that rely on predefined stopping
conditions [15], or fixed time/step limits [16], we employ an
information-theoretic stopping criterion on $\alpha$, which is based
on the ergodic metric. The value of $\alpha$ governs the trade-off
between exploration and exploitation. Larger values of $\alpha$
correspond to an early exploration phase and favour earlier
termination, but may lead to incomplete spatial coverage. Conversely,
smaller values of $\alpha$ indicate that the exploration has largely
converged toward the target EID distribution, resulting in extended
exploitation phases, longer execution times, and only marginal
refinement of already identified regions. In practice, the stopping
threshold on $\alpha$ is selected within the exploitation-dominated
regime, such that further reductions correspond to diminishing
returns in reconstruction accuracy. This choice is supported by
empirical observations indicating saturation of the RMSE as $\alpha$
approaches a plateau, and reflects a practical balance between
coverage completeness, estimation accuracy, and exploration duration.
The advantage of using this metric lies in its ability to quantify
how thoroughly the area has been explored, without relying on prior
knowledge such as fixed time limits or the number and characteristics
of stiff regions that need to be found. This criterion is preferable
to a fixed-time limit, as it adapts the exploration duration to the
complexity of the stiffness distribution and the desired level of map
details, avoiding both premature termination and unnecessary probing.

\vspace{0.8em}

% ---------- Section IV ----------
\begin{center}
\textsc{IV.\ Simulation Results}
\end{center}

\noindent\textit{A.~Setup and Baselines}

\vspace{0.2em}

\noindent
The proposed search algorithm was evaluated on synthetic stiffness
distributions representative of biological soft tissues.

% ---------- Bottom license strip ----------
\vfill
\hrule
\vspace{0.3em}
\begin{center}
{\footnotesize\itshape
This work is licensed under a Creative Commons Attribution 4.0
License. For more information, see
\url{https://creativecommons.org/licenses/by/4.0/}\par}
\end{center}

\end{document}
